\documentclass[10pt]{article}

\usepackage[utf8]{inputenc}

\usepackage[T1]{fontenc}

\usepackage{amsmath}

\usepackage{amsfonts}

\usepackage{amssymb}

\usepackage[version=4]{mhchem}

\usepackage{stmaryrd}



\begin{document}

лова применимы и в общем случае непрерывных отображений $f$ круга $D$, что и было позднее использовано во многих работах.



Точку $\zeta \in \Gamma=\{|z|=1\}$ наз. т о ч к о й $\Pi$ ле с н е р а и относят к множеству $I(f)$, если пересечение



\[

\cap C\left(f, \zeta ; \Delta\left(\zeta, \varphi_{1}, \varphi_{2}\right)\right)

\]



по всем угловым областям $\Delta\left(\zeta, \varphi_{1}, \varphi_{2}\right)$ с вершиной $\zeta$ совпадает с $\Omega$. В 1927 А. И. Плеснер доказал, что для любой мероморфной в круге $D$ функции $f$ почти все точки окружности $\Gamma$ принадлежат либо $F(f)$, либо $I(f)$, то есть $\Gamma=F(f) \bigcup I(f) \cup E$, mes $E=0$. Точку $\zeta \in \Gamma$ наз. то о к о й $\mathrm{M}$ ей $е \mathrm{p} a$ и относят к множеству $M(f)$, если $C(f, \zeta ; D) \neq \Omega$ и пересечение хордальных П. м. $\cap C(f, \zeta ; h(\zeta, \varphi))$ по всем хордам, проведенным в точку $\zeta$, совпадает с $C(f, \zeta ; D)$. К. Мейер (К. Meier, 1961) установил следующий аналог теоремы Плеснера в терминах категории по Бәру: если $f$ мероморфна в $D$, то все точки окружности $\Gamma$, за возможным исключением множества $E$ первой категории, принадлежат объединению $M(f) \cup I(f)$. Получено уточнение теоремы Мейера, согласно к-рому $E$ есть множество первой категории и типа $F_{\sigma}$ (см. [12] - [14], где получены усиления теорем Плеснера и Мейера, а также даны обращение теоремы Мейера и характеризация множества $M(f))$.



Работа П. Фату послужила первоначальным толчком к развитию фундаментальных исследований граничных свойств аналитич. функций. Исследования Ф. и М. Риссов, Н. Н. Л узина, И. И. Привалова, Р. Неванлинны (R. Nevanlinna), А. И. Плеснера, В. И. Смирнова и др. проводились независимо от идей П. Пенлеве, и для них характерно использование методов, связанных с теорией меры и теорией пнтегрирования, с понятием категории по Бәру (см. [4] - [9]).



Основным объектом исследований Ф. Иверсена (F. Iversen) и В. Гpocca (W. Gross) были также мероморфные функции $f$ в области $D$ с жордановой границей $\Gamma=\operatorname{Fr} D$. В произвольной точке $\zeta_{0} \in \Gamma \quad \mathrm{r} p$ а в и чное пг е дельное мно же с т в о $C\left(f, \zeta_{0} ; \Gamma\right)$ определяется следующим образом: если $M_{r}$ обозначает замыкание объединения $\cup C(f, \zeta ; D)$ по всем точкам



\[

\zeta \in\left(\Gamma \backslash\left\{\zeta_{0}\right\}\right) \cup\left\{\left|z-\zeta_{0}\right|<r\right\}

\]



то $C\left(f, \zeta_{0} ; \Gamma\right)=\bigcap_{r>0} M_{r}$. Одна из основных теорем. полученных независимо Ф. Иверсеном и В. Гроссом, утверждает, что при указанных условиях в каждой точке $\zeta_{0} \in \Gamma$ множество



\[

C_{i}\left(f, \zeta_{0} ; D\right)=C\left(f, \zeta_{0} ; D\right) \backslash C\left(f, \zeta_{0} ; \Gamma\right)

\]



открыто и все значения $a \in C_{i}\left(f, \zeta_{0} ; D\right)$, за возможными двумя исключениями, принадлежат множеству повторяющихся значений $R\left(f, \zeta_{0} ; D\right)$; кроме того, каждое исключительное значение (если таковые существуют) является асимптотич. значением функции $f$ в точке $\zeta_{0}$.



Исследования Ф. Иверсена и В. Гросса получили свое дальнейпее развитие в работах А. Бёрлинга (A. Beurling), В. Зайделя (W. Seidel, он и ввел термин «П. м.» в 1932) и др. (см. [5] - [9]). Рассматривались в основном случаи, когда $\zeta_{0}$ принадлежит нек-рому "малому" множеству $E$ точек границы $\Gamma$ нулевой линейной меры или нулевой емкости, и изучалось П. м. $C\left(f, \zeta_{0} ; \Gamma \backslash E\right)$, определяемое аналогично множеству $C(f, \zeta ; \Gamma)$. В этих исследованиях использовались и методы теории потенциала.



Новейшие результаты в этом направлении сформулированы ниже для случая круга $D=\{|z|<1\}$. Пусть фиксировано множество $E$ на дуге $\gamma$ границы $\Gamma=\{|z|=$ $=1\}, \quad$ nes $E=0, \zeta_{0} \in E$. Каждой точке $\zeta \in \gamma \backslash E$ отнесем жорданову дугу $\Lambda_{\zeta} \subset D$, оканчивающуюся в $\zeta$. Пусть $\quad M_{r}^{*}$ - замыкание объединения $\cup C\left(f, \zeta ; \Lambda_{\zeta}\right)$ по всем точкам



\[

\zeta \in(\gamma \backslash E) \cap\left\{\left|z-\zeta_{0}\right|<r\right\}

\]



и пусть $C^{*}\left(f, \zeta_{0}, \Gamma \backslash E\right)=\cap_{r}>0^{M_{r}^{*}}$.



Тогда множество



\[

S\left(\zeta_{0}\right)=C\left(f, \zeta_{0} ; D\right) \backslash C^{*}\left(f, \zeta_{0} ; \Gamma \backslash E\right)

\]



открыто, множество $S\left(\zeta_{0}\right) \backslash R\left(f, \zeta_{0} ; D\right)$ имеет емкость нуль, а каждое значение $a \in S\left(\zeta_{0}\right) \backslash R\left(f, \zeta_{0} ; D\right)$ является асимптотическим значением функции $f$ либо в точке $\zeta_{0}$, либо в каждой точке нек-рой последовательности $\left\{\zeta_{n}\right\}, \zeta_{n} \in \Gamma, n=1,2, \ldots, \lim _{n \rightarrow \infty} \zeta_{n}=\zeta_{0}$. Если емкость $E$ равна нулю, то для каждой связной компоненты $S_{k}\left(\zeta_{0}\right), k=1,2, \ldots$ множества $S\left(\zeta_{0}\right)$ множество $S_{k}\left(\zeta_{0}\right) \backslash R\left(f, \zeta_{0} ; D\right)$ состоит самое большее из двух различных значений.



С помощью нормальных семейств была доказана т е о р е м а Л и нде л ёф а: если голоморфная функция $f$ ограничена в круге $D$ и имеет асимптотич. значение $a$ в точке $\zeta_{0} \in \Gamma$, то она имеет в әтой же точке значение $a$ и в качестве углового граничного значения. Нормальность семейства $F=\{f(z)\}$ мероморфных функций $f(z)$ в области $G$ можно характеризовать в терминах т. в. с ф е р и ч е с к о й пр о и 3 в о дस 0 Й



\[

\rho(f(z))=\frac{\left|f^{\prime}(z)\right|}{1+|f(z)|^{2}} .

\]



Именно, семейство $F$ нормально тогда и только тогда, когда сферич. производные $\rho(f(z)), f \in F$, равномерно ограничены внутри $G$, т. е. для любого компакта $K \subset G$ можно указать такую постоянную $C=C(K)$, что



\[

\rho(f(z)) \leqslant C(K), z \in K, f \in F .

\]



Однако наиболее важный вклад нормальных семейств в теорию П.м. формулируется при помощи понятия нормальной функции. Мероморфная в односвязной области $G$ функция $f(z)$ наз. н о р м а л ьн $о$ й ф у н к ц и е й в $G$, если нормально семейство $\{f(S(z))\}$, где $S$ пробегает семейство всех конформных автоморфизмов области $G ; f(z)$ нормальна в многосвязной области $G$, если она нормальна на универсальной накрывающей поверхности $G$. Мероморфная функция $f(z)$ в круге $D=\{|z|<1\}$ нормальна тогда и только тогда, когда существует константа $C=C(f), 0<C<\infty$, такая, что



\[

\rho(f(z))|d z| \leqslant C \frac{|d z|}{1-|z|^{2}} .

\]



Здесь левая часть есть элемент длины в т. н. хордальной метрике ва сфере Римана $\Omega$ для отображения $w=f(z)$, а стоящее в правой части выражение $d \sigma(z)=$ $=|d z| /\left(1-|z|^{2}\right)$ есть элемент длины в гиперболич. метрике круга $D$. Ограниченные голоморфные функции и мероморфные функции, не принимающие трех различных значений, являются нормальными фуукциями, и нек-рые свойства функций названных классов переносятся на произвольные нормальные функции. Напр., для произвольной нормальной функции справедливо утверждение теоремы Линделёфа. Класс всех нормальных мероморфных функций в круге $D$ имеет нек-рое сходство с классом ограниченного вида функций. Однако имеются и существенные различия. Напр., существуют нормальные мероморфные функции без асимптотич. значений, а следовательно и без радиальных граничных значений, чего не может быть для функций ограниченного вида. Важные исследования асимптотич. значений были проведены Дж. МакЛейном (G. R. MacLane, см. [7], [9]). Теория МакЛейна позволила получить новые доказательства из-





\end{document}